\honest{Quantum Non-Realism}{Eliezer Yudkowsky}{2008}

I open with a question from Einstein, ``asked of Niels Bohr'': ``Does the moon exist when no one is looking at it?'' \nb{The documented source is Abraham Pais, who recalled Einstein asking him, on a walk, ``whether I really believed that the moon exists only when I look at it.'' No source was found for a question put to Bohr.}

Suppose you are working out quantum mechanics in the 1920s, with a theory that gives only frequencies. ``The correct answer is not available to you as a hypothesis, because it will not be invented for another thirty years.'' By the correct answer I mean many-worlds. \nb{Many-worlds was, and is, one interpretation among several, held by a minority of physicists in the polls available.}

The best you could do is a strict ``shut up and calculate'': keep looking for theories, but affirm nothing that the data and the calculations do not force. \nb{The slogan is David Mermin's, from 1989. The stance is close to Bas van Fraassen's constructive empiricism, an anti-realist view.} Then nothing you say can conflict with the correct theory when it arrives. Even saying that the wavefunction ``only gives us probabilities'' went too far, because ``the wavefunction gives us a certainty of many worlds existing.'' \nb{A few lines later the post says only that decoherence ``is not ruled out by the data.'' Its rule is to ``refuse to affirm, as positive knowledge, any proposition which was not forced by the data and the calculations.'' By that rule, ``a certainty of many worlds'' is a claim a strict calculator would not make.}

But it is hard to really shut up and calculate, and people treat their ignorance as knowledge. I show how in dialogues that, I admit, may never have happened. Gallant says, ``I don't know what these equations mean, just that they seem to work.'' Five minutes later Goofus tells a student that the equations ``don't mean anything,'' and, asked how he knows, answers, ``Gallant told me.'' In the same way, a rule for calculating after a measurement becomes the claim that once you know the result, every other probability is zero. I call this ``one of the most embarrassing wrong turns in the history of science.'' \nb{The post names no physicist who made this slide. Without the step that something physical vanishes, the chain is close to QBism, a view developed from 2002, on which a change of quantum state at measurement ``merely reflects the agent's updated epistemic state.'' The post itself says ``partial knowledge is the meaning of probability.''}

Read literally, the chain says consciousness causes collapse. ``Just about nobody'' admits to believing that, but textbooks still talk as if it were so. I name no textbook. Still, I prefer that view to the one that follows, because it at least says the amplitudes, the collapse and the consciousness are real.

The view that follows is Goofus's: ``The equations don't mean anything.'' When a student asks about entangled photons and Bell's theorem, he answers, ``It's meaningless to talk about it before we measure it!'' What state of reality, I ask, would make his words true? And where is he ``getting all this stuff''? Was there an experiment that showed the Schrödinger equation to be meaningless? \nb{The anti-realists had reasons. Bohr treated the wave function as ``a tool for deriving predictions'' because it lives in configuration space and uses imaginary numbers. Pauli compared asking whether something exists that one cannot know anything about to asking ``how many angels are able to sit on the point of a needle.'' QBists answer the question about entangled photons: what changes is the near observer's expectations, and nothing travels to the far photon. All of this is disputed. The post discusses none of it.}

As in ``The Simple Truth,'' ``reality'' is my name for whatever determines my experimental results. So if the equations are ``not real,'' what does determine my results? Before Bell's theorem, the equations could have been partial knowledge of something deterministic, as in statistical mechanics, and Bell's theorem makes that ``much less plausible.'' \nb{Bell's theorem rules out local hidden variables. Bohm's deterministic theory of 1952 survives it as a nonlocal theory, and Bell was its chief advocate.} Either way the equations would mean something. So how can you rule out that they predict so well because they are meaningful? \nb{This is the ``no miracles'' argument for scientific realism, which Hilary Putnam stated in 1975 and philosophers have argued about since.}

``Meaningless'' is a word that stops questions. ``Tell people to `shut up and calculate' because you don't know what the calculations mean, and inside of five years, `Shut up!' will be masquerading as a positive theory of quantum mechanics.'' \nb{In the history the order was reversed: Bohr's reading dates from 1927, the slogan from 1989.} I keep my scorn for those who took ``We don't know why it works'' as knowledge that the equations were ``definitely not real.'' I name no one.

My moral: it is very hard ``to stay in a state of confessed confusion, without making up a story that gives you closure.''
